\nonstopmode{}
\documentclass[a4paper]{book}
\usepackage[times,inconsolata,hyper]{Rd}
\usepackage{makeidx}
\usepackage[utf8]{inputenc} % @SET ENCODING@
% \usepackage{graphicx} % @USE GRAPHICX@
\makeindex{}
\begin{document}
\chapter*{}
\begin{center}
{\textbf{\huge Package `changepointGA'}}
\par\bigskip{\large \today}
\end{center}
\inputencoding{utf8}
\ifthenelse{\boolean{Rd@use@hyper}}{\hypersetup{pdftitle = {changepointGA: Changepoint detection via modified genetic algorithm}}}{}
\begin{description}
\raggedright{}
\item[Type]\AsIs{Package}
\item[Title]\AsIs{Changepoint detection via modified genetic algorithm}
\item[Version]\AsIs{0.1.0}
\item[Date]\AsIs{2024-03-21}
\item[Author]\AsIs{Mo Li }\email{mo.li@louisiana.edu}\AsIs{}
\item[Maintainer]\AsIs{Mo Li }\email{mo.li@louisiana.edu}\AsIs{}
\item[Description]\AsIs{Modified genetic algorithm special for changepoint detection in time series.}
\item[License]\AsIs{MIT + file LICENSE}
\item[RoxygenNote]\AsIs{7.3.1}
\item[Depends]\AsIs{R (>= 4.3.0)}
\item[Imports]\AsIs{foreach, doMC, Rcpp, RcppArmadillo, clue, stats}
\item[LinkingTo]\AsIs{Rcpp, RcppArmadillo}
\item[URL]\AsIs{}\url{https://github.com/mli171/changepointGA}\AsIs{}
\item[BugReports]\AsIs{}\url{https://github.com/mli171/changepointGA/issues}\AsIs{}
\item[Suggests]\AsIs{knitr,
rmarkdown}
\item[Encoding]\AsIs{UTF-8}
\item[LazyData]\AsIs{true}
\item[VignetteBuilder]\AsIs{knitr}
\end{description}
\Rdcontents{\R{} topics documented:}
\inputencoding{utf8}
\HeaderA{AMOCcrossover}{Average crossover operator to produce offspring for AMOC problem}{AMOCcrossover}
%
\begin{Description}
In this crossover operator designed for AMOC problem, the new child is
produced by taking the average of the changepoint locations from \code{dad}
and \code{mom} and round to an integer. Note, every chromosome has at most
one candidate changepoint location.
\end{Description}
%
\begin{Usage}
\begin{verbatim}
AMOCcrossover(mom, dad, p.range = NULL, minDist, lmax, N)
\end{verbatim}
\end{Usage}
%
\begin{Arguments}
\begin{ldescription}
\item[\code{mom}] Among two selected individuals, \code{mom} represents the selected
chromosome representation with lower fitness function value.

\item[\code{dad}] Among two selected individuals, \code{dad} represents the selected
chromosome representation with larger fitness function value.

\item[\code{p.range}] The default value is \code{NULL}. If there is no requirement
on model order selection, such an auxiliary argument is needed for \code{GA}
and \code{IslandGA} functions.

\item[\code{minDist}] The minimum length between two adjacent changepoints.

\item[\code{lmax}] The maximum possible length of the chromosome representation.

\item[\code{N}] The length of time series.
\end{ldescription}
\end{Arguments}
%
\begin{Value}
The child chromosome that produced from \code{mom} and \code{dad} for
next generation.
\end{Value}
\inputencoding{utf8}
\HeaderA{AMOCmutation}{Jump mutation operator to produce offspring for AMOC problem}{AMOCmutation}
%
\begin{Description}
In a certain probability, the \code{mutation} genetic operator can be applied
to generate a new \code{child}. In this AMOC mutation operator, the new child
changepoint location can be down via a "jump" method. The child changepoint
location will jump \code{minDist} time units either to the left or right to
produce the changepoint location for the mutated child. The jump direction is
randomly decided with 0.5 probability.
\end{Description}
%
\begin{Usage}
\begin{verbatim}
AMOCmutation(
  child,
  p.range = NULL,
  minDist,
  Pchangepoint = NULL,
  lmax = NULL,
  mmax = NULL,
  N = NULL
)
\end{verbatim}
\end{Usage}
%
\begin{Arguments}
\begin{ldescription}
\item[\code{child}] The child chromosome resulting from the \code{crossover} genetic
operator.

\item[\code{p.range}] The default value is \code{NULL}. If there is no requirement
on model order selection, such an auxiliary argument is needed for \code{GA}
and \code{IslandGA} functions.

\item[\code{minDist}] The minimum length between two adjacent changepoints in
\code{\LinkA{AMOCselection}{AMOCselection}} operator, which is also the jump magnitude in the
\code{AMOCmutation} operator.

\item[\code{Pchangepoint}] An auxiliary argument is needed for \code{GA}
and \code{IslandGA} functions.

\item[\code{lmax}] An auxiliary argument is needed for \code{GA} and \code{IslandGA}
functions.

\item[\code{mmax}] An auxiliary argument is needed for \code{GA} and \code{IslandGA}
functions.

\item[\code{N}] An auxiliary argument is needed for \code{GA} and \code{IslandGA}
functions.
\end{ldescription}
\end{Arguments}
%
\begin{Value}
The resulting child chromosome representation.
\end{Value}
\inputencoding{utf8}
\HeaderA{AMOCpopulation}{Random population initialization for AMOC problem}{AMOCpopulation}
%
\begin{Description}
Randomly generate the individuals' chromosomes (changepoint configurations)
to construct the first generation population for the at most one changepoint
(AMOC) problem.
\end{Description}
%
\begin{Usage}
\begin{verbatim}
AMOCpopulation(popsize, p.range, N, minDist, Pchangepoint, mmax, lmax)
\end{verbatim}
\end{Usage}
%
\begin{Arguments}
\begin{ldescription}
\item[\code{popsize}] An integer represents the number of individual in each
population for GA (or subpopulation for IslandGA).

\item[\code{p.range}] Default is \code{NULL} for only changepoint detection. If
\code{p.range} is specified as a list object, which contains the range of
each model order parameters for order selection (integers). The number of
order parameters must be equal to the length of \code{p.range}.

\item[\code{N}] The length of time series.

\item[\code{minDist}] The minimum length between two adjacent changepoints.

\item[\code{Pchangepoint}] The probability that a changepoint can occur.

\item[\code{mmax}] The maximum possible number of changepoints in the data set.

\item[\code{lmax}] The maximum possible length of the chromosome representation.
\end{ldescription}
\end{Arguments}
%
\begin{Details}
Given the possible candidate changepoint location set, each chromosome in the
first generation population can be obtained by randomly sampling one location
from the candidate set. The first element of every chromosome represent the
number of changepoints and the last non-zero element always equal to the
length of time series plus one (N+1).
\end{Details}
%
\begin{Value}
A matrix that contains each individual's chromosome.
\end{Value}
\inputencoding{utf8}
\HeaderA{AMOCselection}{The parents selection genetic algorithm operator for AMOC problem}{AMOCselection}
%
\begin{Description}
The genetic algorithm require to select a pair of chromosomes, representing
\code{dad} and \code{mom}, for the \code{crossover} operator to
produce offspring (individual for next generation). Here, the same linear
ranking method in \code{\LinkA{selection\_linearrank}{selection.Rul.linearrank}} is used to select a pair
of chromosomes for \code{dad} and \code{mom} in the at most one changepoint
(AMOC) problem. By default, the dad has better fit/smaller fitness function
value/larger rank than \code{mom}.
\end{Description}
%
\begin{Usage}
\begin{verbatim}
AMOCselection(pop, popFit)
\end{verbatim}
\end{Usage}
%
\begin{Arguments}
\begin{ldescription}
\item[\code{pop}] A matrix contains the chromosomes for all individuals. The number of
rows is equal to \code{lmax} and the number of columns is equal to the
\code{popsize}.

\item[\code{popFit}] A vector contains the objective function value (population fit)
being associated to each individual chromosome from above.
\end{ldescription}
\end{Arguments}
%
\begin{Value}
A list contains the chromosomes for \code{dad} and \code{mom}.
\end{Value}
\inputencoding{utf8}
\HeaderA{ARIMA.BIC.Order}{Calculating BIC for Multiple changepoint detection with model order selection}{ARIMA.BIC.Order}
%
\begin{Description}
The objective function for changepoint search in Autoregressive
moving average with model order selection via Bayesian Information Criterion
(BIC).
\end{Description}
%
\begin{Usage}
\begin{verbatim}
ARIMA.BIC.Order(chromosome, plen = 2, XMat, Xt)
\end{verbatim}
\end{Usage}
%
\begin{Arguments}
\begin{ldescription}
\item[\code{chromosome}] The chromosome consists of the number of changepoints, the
order of AR part, the order of MA part, the changepoint locations, and a
value of time series length plus 1 (N+1) indicating the end of the chromosome.

\item[\code{plen}] The number of model order parameters that need to be selected.
If model order selection needs to be performed simultaneously with the
changepoint detection task, \code{plen} should be nonzero.

\item[\code{XMat}] A matrix contains the covariates, but not includes changepoint
effects, for time series regression.

\item[\code{Xt}] The simulated ARMA time series from \code{ts.sim} function.
\end{ldescription}
\end{Arguments}
%
\begin{Value}
The BIC value of the objective function.
\end{Value}
%
\begin{Examples}
\begin{ExampleCode}
N = 1000
XMatT = matrix(1, nrow=N, ncol=1)
Xt = ts.sim(beta=0.5, XMat=XMatT, sigma=1, phi=0.5, theta=0.8,
            Delta=c(2, -2), CpLoc=c(250, 750), seed=1234)

# one chromosome representation
chromosome = c(2, 1, 1, 250, 750, 1001)
ARIMA.BIC.Order(chromosome, plen=2, XMat=XMatT, Xt=Xt)
\end{ExampleCode}
\end{Examples}
\inputencoding{utf8}
\HeaderA{arrayaCGH}{Glioblastoma array aCGH data}{arrayaCGH}
\keyword{datasets}{arrayaCGH}
%
\begin{Description}
Glioblastoma multiforme (GBM) can cause chromosomal aberrations during DNA
replication. Array comparative genomic hybridization (aCGH) experiments analyze
these aberrations using fluorescence intensity ratios, which indicate DNA copy
number variations.
\end{Description}
%
\begin{Usage}
\begin{verbatim}
data(arrayaCGH)
\end{verbatim}
\end{Usage}
%
\begin{Format}
An object of class \code{data.frame} with 193 rows and 5 columns.
\end{Format}
%
\begin{References}
Lai, W. R., Johnson, M. D., Kucherlapati, R., \& Park, P. J. (2005).
Comparative analysis of algorithms for identifying amplifications and deletions
in array CGH data. \emph{Bioinformatics}, 21(19), 3763-3770.

Killick, R., \& Eckley, I. A. (2014). changepoint: An R package for
changepoint analysis. \emph{Journal of statistical software}, 58, 1-19.
\end{References}
%
\begin{Examples}
\begin{ExampleCode}
data(arrayaCGH)
\end{ExampleCode}
\end{Examples}
\inputencoding{utf8}
\HeaderA{AtlantaTemperature}{Atlanta airport yearly average temperatures data}{AtlantaTemperature}
\keyword{datasets}{AtlantaTemperature}
%
\begin{Description}
The annual mean temperature data from 1879 to 2013 at Atlanta, Georgia’s
Hartsfield International Airport weather station.
\end{Description}
%
\begin{Usage}
\begin{verbatim}
data(AtlantaTemperature)
\end{verbatim}
\end{Usage}
%
\begin{Format}
An object of class \code{data.frame} with 134 rows and 2 columns.
\end{Format}
%
\begin{References}
Lund, R. B., Beaulieu, C., Killick, R., Lu, Q., \& Shi, X. (2023).
Good practices and common pitfalls in climate time series changepoint techniques:
A review. \emph{Journal of Climate}, 36(23), 8041-8057.
\end{References}
%
\begin{Examples}
\begin{ExampleCode}
data(AtlantaTemperature)
\end{ExampleCode}
\end{Examples}
\inputencoding{utf8}
\HeaderA{BinSearch.BIC}{Example function: Calculating BIC for AR(1) model}{BinSearch.BIC}
%
\begin{Description}
The example objective function for changepoint search via Bayesian
Information Criterion (BIC) for simple AR(1) model from Shi et al. (2022).
The model is expressed as \eqn{X_{t}=\kappa_{t}+\epsilon_{t}}{}, where
\eqn{\epsilon_{t}}{} follows a stationary AR(1) process, and \eqn{\kappa_{t}}{}
denotes the regimen mean.
\end{Description}
%
\begin{Usage}
\begin{verbatim}
BinSearch.BIC(chromosome, plen = 0, Xt)
\end{verbatim}
\end{Usage}
%
\begin{Arguments}
\begin{ldescription}
\item[\code{chromosome}] The chromosome consists of the number of changepoints, their
locations, and a value of time series length plus 1 (N+1) indicating the end
of the chromosome.

\item[\code{plen}] The number of model order parameters that need to be selected.
Since we don't need model order selection in this example, \code{plen} equals
to 0.

\item[\code{Xt}] The simulated AR(1) time series from \code{ts.sim} function.
\end{ldescription}
\end{Arguments}
%
\begin{Value}
The BIC value of the objective function.
\end{Value}
%
\begin{References}

Shi, X., Gallagher, C., Lund, R., \& Killick, R. (2022).
A comparison of single and multiple changepoint techniques for time series data.
\emph{Computational Statistics \& Data Analysis}, 170, 107433.

\end{References}
%
\begin{Examples}
\begin{ExampleCode}
Ts = 1000
betaT = c(0.5) # intercept
XMatT = matrix(1, nrow=Ts, ncol=1)
colnames(XMatT) = "intercept"
sigmaT = 1
phiT = c(0.5)
thetaT = NULL
DeltaT = c(2, -2)
Cp.prop = c(1/4, 3/4)
CpLocT = floor(Ts*Cp.prop)

myts = ts.sim(beta=betaT, XMat=XMatT, sigma=sigmaT, phi=phiT, theta=thetaT,
              Delta=DeltaT, CpLoc=CpLocT, seed=1234)

# candidate changepoint configuration
chromosome = c(2, 250, 750, 1001)
BinSearch.BIC(chromosome, Xt=myts)
\end{ExampleCode}
\end{Examples}
\inputencoding{utf8}
\HeaderA{cptDist}{Comparing multiple changepoint configurations by pairwise distance}{cptDist}
%
\begin{Description}
This function is to calculate the distance metric (Shi et al., 2022) between
two changepoint configurations, \eqn{C_{1}=\{\tau_{1}, \ldots, \tau_{m}\}}{}
and \eqn{C_{2}=\{\eta_{1}, \ldots, \eta_{k}\}}{}, where \eqn{\tau}{}'s are the
changepoint locations.
\end{Description}
%
\begin{Usage}
\begin{verbatim}
cptDist(tau1, tau2, N)
\end{verbatim}
\end{Usage}
%
\begin{Arguments}
\begin{ldescription}
\item[\code{tau1}] A vector contains the changepoint locations for \eqn{C_{1}}{}
for comparison. A value \code{NULL} is required if there is no changepoint detected.

\item[\code{tau2}] A vector contains the changepoint locations for  \eqn{C_{2}}{}
for comparison. A value \code{NULL} is required if there is no changepoint detected.

\item[\code{N}] The simulated time series sample size. Two changepoint configurations
should have the same \eqn{n}{} values.
\end{ldescription}
\end{Arguments}
%
\begin{Details}
The pairwise distance was proposed by Shi et al. (2022),
\deqn{d(C_{1}, C_{2})=|m-k|+ \min(A(C_{1}, C_{2})),}{}
where \eqn{m}{} is the number of changepoints in configuration \eqn{C_{1}}{} and
\eqn{k}{} is the number of changepoints in configuration \eqn{C_{2}}{}.
The term \eqn{\min(A(C_{1}, C_{2}))}{} reflects the cost of matching
changepoint locations between \eqn{C_{1}}{} and \eqn{C_{2}}{} and can be calculated
using the linear assignment method. Details can be found in Shi et al. (2022).
Note: if one configuration doesn't contain any changepoints (valued \code{NULL}),
the distance is defined as \eqn{|m-k|}{}. The function can also be used
to examine changepoint detection performance in simulation studies. Given the
true changepoint configuration, \eqn{C_{true}}{}, used in generating the time
series, the calculated distance between the estimated multiple changepoint
configuration, \eqn{C_{est}=\{\eta_{1}, \ldots, \eta_{k}\}}{}, and \eqn{C_{true}}{}
can be used to evaluate the performance of the detection algorithm.
\end{Details}
%
\begin{Value}
\begin{ldescription}
\item[\code{dist}] The calculated distance.
\end{ldescription}
\end{Value}
%
\begin{References}

Shi, X., Gallagher, C., Lund, R., \& Killick, R. (2022).
A comparison of single and multiple changepoint techniques for time series data.
\emph{Computational Statistics \& Data Analysis}, 170, 107433.

\end{References}
%
\begin{Examples}
\begin{ExampleCode}
N = 100

# both tau1 and tau2 has detected changepoints
tau2 = c(25, 50, 75)
tau1 = c(20, 35, 70, 80, 90)
cptDist(tau1=tau1, tau2=tau2, N=N)

# either tau1 or tau2 has zero detected changepoints
cptDist(tau1=tau1, tau2=NULL, N=N)
cptDist(tau1=NULL, tau2=tau2, N=N)

# both tau1 and tau2 has zero detected changepoints
cptDist(tau1=NULL, tau2=NULL, N=N)
\end{ExampleCode}
\end{Examples}
\inputencoding{utf8}
\HeaderA{GA}{Genetic algorithm}{GA}
%
\begin{Description}
Perform the modified genetic algorithm for changepoint detection.
This involves the minimization of an objective function using a genetic algorithm (GA).
The algorithm can be run sequentially or with explicit parallelization.
\end{Description}
%
\begin{Usage}
\begin{verbatim}
GA(
  ObjFunc,
  N,
  GA_param = .default.GA_param,
  GA_operators = .default.operators,
  p.range = NULL,
  ...
)
\end{verbatim}
\end{Usage}
%
\begin{Arguments}
\begin{ldescription}
\item[\code{ObjFunc}] The fitness function to be minimized. Users can specify any R or Rcpp
function as the fitness function, setting the input as the potential solution to
the optimization problem and returning a numerical value as the output/fitness.
Depending on the user-specified chromosome representation, the optimization task
can be changepoint detection only or changepoint detection plus model order selection,
which can be specified via the \code{option} parameter in \code{\LinkA{GA\_param}{GA.Rul.param}}. When
\code{option="both"}, the list \code{p.range} must be specified to give the range
of model orders.

\item[\code{N}] The sample size of the time series.

\item[\code{GA\_param}] A list contains the hyper-parameters for genetic algorithm.
The default values for these hyper-parameters are included in \code{.default.GA\_param}.
See \code{\LinkA{GA\_param}{GA.Rul.param}} for the details.

\item[\code{GA\_operators}] A list includes the functions for population initialization,
new individual selection, and genetic operator of crossover and mutation.
See \code{\LinkA{operators}{operators}} for the details and default functions.

\item[\code{p.range}] Default is \code{NULL} for only changepoint detection. If
\code{p.range} is specified as a list object, which contains the range of
each model order parameters for order selection (integers). The number of
order parameters must be equal to the length of \code{p.range}.

\item[\code{...}] additional arguments that will be passed to the fitness function.
\end{ldescription}
\end{Arguments}
%
\begin{Value}
Returns a list that has components:
\begin{ldescription}
\item[\code{overbestfit}] The obtained minimum value of the objective function after
the final iteration.
\item[\code{overbestchrom}] The locations of the detected changepoints associating
with the \code{overbestfit} the after the final iteration.
\item[\code{bestfit}] The minimized fitness function values at each iteration.
\item[\code{bestchrom}] The detected changepoints at each iteration.
\item[\code{count}] The number of iterations undertaken by the genetic algorithm.
\item[\code{convg}] An integer code.
\begin{itemize}

\item{} 0 indicates the algorithm successful completion.
\item{} 1 indicates the the total number of generations exceeds the
prespecified \code{maxgen} limit.

\end{itemize}


\end{ldescription}
\end{Value}
%
\begin{Examples}
\begin{ExampleCode}
N = 1000
XMatT = matrix(1, nrow=N, ncol=1)
Xt = ts.sim(beta=0.5, XMat=XMatT, sigma=1, phi=0.5, theta=NULL,
            Delta=c(2, -2), CpLoc=c(250, 750), seed=1234)
TsPlotCheck(X=1:N, Xat=seq(from=1, to=N, length=10), Y=Xt, tau=c(250, 750))

GA.res = GA(ObjFunc=BinSearch.BIC, N=N, Xt=Xt)
GA.res$overbestfit
GA.res$overbestchrom
\end{ExampleCode}
\end{Examples}
\inputencoding{utf8}
\HeaderA{GA\_param}{GA\_param}{GA.Rul.param}
%
\begin{Description}
A list object contains the hyperparameters for GA running.
\end{Description}
%
\begin{Arguments}
\begin{ldescription}
\item[\code{popsize}] An integer represents the number of individuals in each population.
\item[\code{Pcrossover}] The probability that the crossover operator applies on two individual chromosomes.
\item[\code{Pmutation}] The probability that the mutation operator applies on one individual chromosome.
\item[\code{Pchangepoint}] The probability that a changepoint has occurred. User
could change this probability based on domain knowledge and the time series
length.
\item[\code{minDist}] The minimum length between two adjacent changepoints.
\item[\code{mmax}] The maximum possible number of changepoints in the data set.
For a time series of length 1000, the default value is 499. The suggested
value should be based on the length of the time series. For instance, if a
time series has length N, the recommended \code{mmax} should be N/2-1.
\item[\code{lmax}] The maximum possible length of the chromosome representation. For
a time series of length 1000, the default value is 501. The suggested value
should be based on the length of the time series. For instance, if a time
series has length N, the recommended \code{lmax} should be 2+N/2-1.
\item[\code{maxgen}] The maximum number of generation that the GA can last.
\item[\code{maxconv}] If the overall best fitted value doesn't change after \code{maxconv} consecutive migrations, the GA algorithm stops.
\item[\code{option}] A string controls the optimization task. ``cp'' indicates the task is changepoint detection only. ``both'' indicates the task will include both changepoint detection and model order selection.
\item[\code{monitoring}] A binary interger 0 or 1, indicating whether print out middle results for each iterations of GA.
\item[\code{parallel}] Whether use multiple threads to parallel compute the individual fittness function values.
\item[\code{nCore}] An integer represents the number of cores used in parallel computing.
\item[\code{tol}] The tolerance level for deciding GA to stop.
\item[\code{seed}] An single integer allows function produce reproducible results.
\end{ldescription}
\end{Arguments}
%
\begin{Author}
Mo Li
\end{Author}
%
\begin{Examples}
\begin{ExampleCode}
# time series length
N = 1000

GA_param = list(
  popsize      = 40,
  Pcrossover   = 0.95,
  Pmutation    = 0.1,
  Pchangepoint = 0.06,
  minDist      = 3,
  mmax         = N/2 - 1,
  lmax         = 2 + N/2 - 1,
  maxgen       = 50000,
  maxconv      = 5000,
  option       = "cp",
  monitoring   = FALSE,
  parallel     = FALSE,
  nCore        = NULL,
  tol          = 1e-5,
  seed         = NULL
)
\end{ExampleCode}
\end{Examples}
\inputencoding{utf8}
\HeaderA{IslandGA}{Island model based genetic algorithm}{IslandGA}
%
\begin{Description}
Perform the modified island-based genetic algorithm (IslandGA) for multiple changepoint detection.
Minimization of an objective function using genetic algorithm (GA).
The algorithm can be run sequentially or in explicit parallelisation.
\end{Description}
%
\begin{Usage}
\begin{verbatim}
IslandGA(
  ObjFunc,
  N,
  IslandGA_param = .default.IslandGA_param,
  IslandGA_operators = .default.operators,
  p.range = NULL,
  ...
)
\end{verbatim}
\end{Usage}
%
\begin{Arguments}
\begin{ldescription}
\item[\code{ObjFunc}] The fitness function to be minimized. Users can specify any R or Rcpp
functions as the fitness function with setting input as potential solution to
the optimization problem and returning a numerical value as the output/fitness.

\item[\code{N}] The sample size of the time series.

\item[\code{IslandGA\_param}] A list contains the hyper-parameters for IslandGA.
The default values for these hyper-parameters are included in \code{.default.IslandGA\_param}.
See \code{\LinkA{IslandGA\_param}{IslandGA.Rul.param}} for the details.

\item[\code{IslandGA\_operators}] A list includes the functions for population initialization,
new individual selection, and genetic operator of crossover and mutation.
See \code{\LinkA{operators}{operators}} for the details and default functions.

\item[\code{p.range}] Default is \code{NULL} for only changepoint detection. If
\code{p.range} is specified as a list object, which contains the range of
each model order parameters for order selection (integers). The number of
order parameters must be equal to the length of \code{p.range}.

\item[\code{...}] additional arguments that will be passed to the fitness function.
\end{ldescription}
\end{Arguments}
%
\begin{Value}
Returns a list that has the following components.
\begin{ldescription}
\item[\code{overbestfit}] The obtained minimum value of the objective function after
the final iteration.
\item[\code{overbestchrom}] The locations of the detected changepoints associating
with the \code{overbestfit} the after the final iteration.
\item[\code{bestfit}] The minimized fitness function values at each iteration.
\item[\code{bestchrom}] The detected changepoints at each iteration.
\item[\code{countMig}] The number of migrations undertaken by the IslandGA.
\item[\code{count}] The number of iterations (generations) undertaken by the island
genetic algorithm model.
\item[\code{convg}] An integer code.
\begin{itemize}

\item{} 0 indicates the algorithm successful completion.
\item{} 1 indicates the the total number of generations exceeds the
prespecified \code{maxgen} limit.

\end{itemize}


\end{ldescription}
\end{Value}
%
\begin{Examples}
\begin{ExampleCode}
N = 1000
XMatT = matrix(1, nrow=N, ncol=1)
Xt = ts.sim(beta=0.5, XMat=XMatT, sigma=1, phi=0.5, theta=NULL,
            Delta=c(2, -2), CpLoc=c(250, 750), seed=1234)
TsPlotCheck(X=1:N, Xat=seq(from=1, to=N, length=10), Y=Xt, tau=c(250, 750))

IslandGA.res = IslandGA(ObjFunc=BinSearch.BIC, N=N, Xt=Xt)
IslandGA.res$overbestfit
IslandGA.res$overbestchrom
\end{ExampleCode}
\end{Examples}
\inputencoding{utf8}
\HeaderA{IslandGA\_param}{IslandGA\_param}{IslandGA.Rul.param}
%
\begin{Description}
A list object contains the hyperparameters for island-based GA running.
\end{Description}
%
\begin{Arguments}
\begin{ldescription}
\item[\code{subpopsize}] An integer represents the number of individual in each sub-population (island).
\item[\code{Islandsize}] The number of subpopulation (island).
\item[\code{Pcrossover}] The probability that the crossover operator applies on two individual chromosomes.
\item[\code{Pmutation}] The probability that the mutation operator applies on one individual chromosome.
\item[\code{Pchangepoint}] The probability that a changepoint has occurred.
\item[\code{minDist}] The minimum length between two adjacent changepoints.
\item[\code{mmax}] The maximum possible number of changepoints in the data set.
\item[\code{lmax}] The maximum possible length of the chromosome representation.
\item[\code{maxMig}] The maximum number of migrations. After \code{maxMig} migrations, the island-based GA algorithm stops.
\item[\code{maxgen}] The maximum number of generations that each subpopulation (island) has. The migration will apply after \code{maxgen} generations for each subpopulation(island).
\item[\code{maxconv}] If the overall best fitted value doesn't change after  \code{maxconv} consecutive migrations, the island-based GA algorithm stops.
\item[\code{option}] A string controls the optimization task. ``cp'' indicates the task is changepoint detection only. ``both'' indicates the task will include both changepoint detection and model order selection.
\item[\code{monitoring}] A binary interger 0 or 1, indicating whether print out middle results for each iterations of GA.
\item[\code{parallel}] Whether use multiple threads to parallel compute the individual fittness function values.
\item[\code{nCore}] An integer represents the number of cores used in parallel computing.
\item[\code{tol}] The tolerance level for deciding GA to stop.
\item[\code{seed}] An single integer allows function produce reproducible results.
\end{ldescription}
\end{Arguments}
%
\begin{Author}
Mo Li
\end{Author}
%
\begin{Examples}
\begin{ExampleCode}
# time series length
N = 1000

IslandGA_param = list(
  subpopsize   = 40,
  Islandsize   = 5,
  Pcrossover   = 0.95,
  Pmutation    = 0.15,
  Pchangepoint = 0.1,
  minDist      = 2,
  mmax         = N/2 - 1,
  lmax         = 2 + N/2 - 1,
  maxMig       = 1000,
  maxgen       = 50,
  maxconv      = 100,
  option       = "cp",
  monitoring   = FALSE,
  parallel     = FALSE, ###
  nCore        = NULL,
  tol          = 1e-5,
  seed         = NULL
)
\end{ExampleCode}
\end{Examples}
\inputencoding{utf8}
\HeaderA{mutation}{The default mutation operator in genetic algorithm}{mutation}
%
\begin{Description}
In a certain probability, the \code{mutation} genetic operator can be applied
to generate a new \code{child}. By default, the new child selection can be
down by the similar individual selection method in population initialization,
\code{\LinkA{selectTau}{selectTau}}.
\end{Description}
%
\begin{Usage}
\begin{verbatim}
mutation(child, p.range = NULL, minDist, Pb, lmax, mmax, N)
\end{verbatim}
\end{Usage}
%
\begin{Arguments}
\begin{ldescription}
\item[\code{child}] The child chromosome resulting from the \code{crossover} genetic
operator.

\item[\code{p.range}] Default is \code{NULL} for only changepoint detection. If
\code{p.range} is specified as a list object, which contains the range of
each model order parameters for order selection (integers). The number of
order parameters must be equal to the length of \code{p.range}.

\item[\code{minDist}] The required minimum distance between two adjacent changepoints.

\item[\code{Pb}] The probability of changepoints for every time series.

\item[\code{lmax}] The user specified maximum number of changepoints, by default,
as \code{N/2 - 1}.

\item[\code{mmax}] The user specified maximum length of individual chromosome,
by default, as \code{2+N/2 + 1}.

\item[\code{N}] The sample size of the time series.
\end{ldescription}
\end{Arguments}
%
\begin{Details}
A function can apply mutation to the produced child with the specified
probability \code{Pmutation} in \code{\LinkA{GA\_param}{GA.Rul.param}} and
\code{\LinkA{IslandGA\_param}{IslandGA.Rul.param}}. If order selection is not requested
(\code{option = "cp"} in \code{GA\_param} and \code{Island\_GA}), the default
\code{\LinkA{mutation}{mutation}} operator function uses \code{selectTau} to select
a completely new individual with a new chromosome as the mutated child.
For details, see \code{\LinkA{selectTau}{selectTau}}. If order selection is needed
(\code{option = "both"} in \code{GA\_param} and \code{Island\_GA}), we first
decide whether to keep the produced child's model order with a probability
of 0.5. If the child's model order is retained, the \code{selectTau}
function is used to select a completely new individual with a new chromosome
as the mutated child. If a new model order is selected from the candidate
model order set, there is a 0.5 probability to either select a completely new
individual with new changepoint locations or retain the original child's
changepoint locations for the mutated child. Note that the current model
orders in the child's chromosome are excluded from the set to avoid redundant
objective function evaluation. Finally, the function returns a vector
containing the modified chromosomes for mutated \code{child}.
\end{Details}
%
\begin{Value}
The resulting child chromosome representation.
\end{Value}
\inputencoding{utf8}
\HeaderA{operators}{operators}{operators}
%
\begin{Description}
A list object contains the GA or IslandGA operators function names.
\end{Description}
%
\begin{Arguments}
\begin{ldescription}
\item[\code{population}] It could be a function or a matrix. The function should be designed for initializing a population. The default population initialization is random initialization with some imposed constraints. See \code{\LinkA{random\_population}{random.Rul.population}} for example. The function returned object is a matrix, \code{pop}. The users can specified their own \code{population} function. It could also be a matrix object, which contain the user specified chromosome. By default, each column represents one individual chromosome.
\item[\code{selection}] A function can help select \code{mom} and \code{dad} from current generation population, where \code{dad} is set to have better fit (smaller fitness function values). The default for selection uses the linear rank selection method. See \code{\LinkA{selection\_linearrank}{selection.Rul.linearrank}} for example. The function returned object is a list contain the chromosomes for \code{mom} and \code{dad}.
\item[\code{crossover}] A function can apply crossover to the chosen parents to produce child for next generation with specified probability. The default for crossover uses the uniform crossover method. See \code{\LinkA{uniformcrossover}{uniformcrossover}} for details in the default crossover operator. The function returned object is a vector contain the chromosomes for \code{child}.
\item[\code{mutation}] A function can apply mutation to the produced child with the specified probability \code{Pmutation} in \code{\LinkA{GA\_param}{GA.Rul.param}} and \code{\LinkA{IslandGA\_param}{IslandGA.Rul.param}}. See \code{\LinkA{mutation}{mutation}} for details in the default mutation operator.
\end{ldescription}
\end{Arguments}
%
\begin{References}
Lu, Q., Lund, R., \& Lee, T. C. (2010). An MDL approach to the climate segmentation problem.
\emph{Ann. Appl. Stat.} 4 (1) 299 - 319.
\end{References}
%
\begin{SeeAlso}
See Also as \code{\LinkA{GA}{GA}}.
\end{SeeAlso}
%
\begin{Examples}
\begin{ExampleCode}
operators = list(population = "random_population",
                 selection  = "selection_linearrank",
                 crossover  = "uniformcrossover",
                 mutation   = "mutation")
\end{ExampleCode}
\end{Examples}
\inputencoding{utf8}
\HeaderA{random\_population}{Random population initialization}{random.Rul.population}
%
\begin{Description}
Randomly generate the individuals' chromosomes (changepoint confirgurations)
to construct the first generation population.
\end{Description}
%
\begin{Usage}
\begin{verbatim}
random_population(popsize, prange, N, minDist, Pb, mmax, lmax)
\end{verbatim}
\end{Usage}
%
\begin{Arguments}
\begin{ldescription}
\item[\code{popsize}] An integer represents the number of individual in each
population for GA (or subpopulation for IslandGA).

\item[\code{prange}] Default is \code{NULL} for only changepoint detection. If
\code{prange} is specified as a list object, which contains the range of
each model order parameters for order selection (integers). The number of
order parameters must be equal to the length of \code{prange}.

\item[\code{N}] The length of time series.

\item[\code{minDist}] The minimum length between two adjacent changepoints.

\item[\code{Pb}] Same as \code{Pchangepoint}, the probability that a changepoint has occurred.

\item[\code{mmax}] The maximum possible number of changepoints in the data set.

\item[\code{lmax}] The maximum possible length of the chromosome representation.
\end{ldescription}
\end{Arguments}
%
\begin{Details}
The default population initialization uses \code{\LinkA{selectTau}{selectTau}} to
select the chromosome for the first generation population. Each column from
the produced population matrix represent an chromosome of an individual.
The first element of every chromosome represent the number of changepoints
and the last non-zero element always equal to the length of time series
plus one (N+1).
\end{Details}
%
\begin{Value}
A matrix that contains each individual's chromosome.
\end{Value}
\inputencoding{utf8}
\HeaderA{selection\_linearrank}{The default parents selection genetic algorithm operator}{selection.Rul.linearrank}
%
\begin{Description}
The genetic algorithm require to select a pair of chromosomes, representing
\code{dad} and \code{mom}, for the \code{crossover} operator to
produce offspring (individual for next generation). The parents chromosomes
are randomly selectd from the initialized population by a linear ranking
method according to each individual's fittness in the input argument
\code{popFit}. By default, the dad has better fit/smaller fitness function
value/larger rank than \code{mom}.
\end{Description}
%
\begin{Usage}
\begin{verbatim}
selection_linearrank(pop, popFit)
\end{verbatim}
\end{Usage}
%
\begin{Arguments}
\begin{ldescription}
\item[\code{pop}] A matrix contains the chromosomes for all individuals. The number of
rows is equal to \code{lmax} and the number of columns is equal to the
\code{popsize}.

\item[\code{popFit}] A vector contains the objective function value (population fit)
being associated to each individual chromosome from above.
\end{ldescription}
\end{Arguments}
%
\begin{Value}
A list contains the chromosomes for \code{dad} and \code{mom}.
\end{Value}
\inputencoding{utf8}
\HeaderA{selectTau}{Randomly select the chromosome}{selectTau}
%
\begin{Description}
Randomly select the changepoint configuration for population initialization.
The selected changepoint configuration represents a changepoint chromosome.
The first element of the chromosome represent the number of changepoints
and the last non-zero element always equal to the length of time series + 1.
\end{Description}
%
\begin{Usage}
\begin{verbatim}
selectTau(N, prange, minDist, Pb, mmax, lmax)
\end{verbatim}
\end{Usage}
%
\begin{Arguments}
\begin{ldescription}
\item[\code{N}] The length of time series.

\item[\code{prange}] A list object containing the possible range for other
pre-defined model parameters, i.e. AR/MA order of ARMA models.

\item[\code{minDist}] The minimum length between two adjacent changepoints.

\item[\code{Pb}] Same as \code{Pchangepoint}, the probability that a changepoint has occurred.

\item[\code{mmax}] The maximum possible number of changepoints in the data set.

\item[\code{lmax}] The maximum possible length of the chromosome representation.
\end{ldescription}
\end{Arguments}
%
\begin{Value}
A single changepoint configuration format as above.
\end{Value}
\inputencoding{utf8}
\HeaderA{ts.sim}{Time series simulation with changepoint effects}{ts.sim}
%
\begin{Description}
This is a function to simulate time series with changepoint effects. See
details below.
\end{Description}
%
\begin{Usage}
\begin{verbatim}
ts.sim(
  beta,
  XMat,
  sigma,
  phi = NULL,
  theta = NULL,
  Delta = NULL,
  CpLoc = NULL,
  seed = NULL
)
\end{verbatim}
\end{Usage}
%
\begin{Arguments}
\begin{ldescription}
\item[\code{beta}] A parameter vector contains other mean function parameters without changepoint parameters.

\item[\code{XMat}] The covairates for time series mean function without changepoint indicators.

\item[\code{sigma}] The standard deviation for time series residuals \eqn{\epsilon_{t}}{}.

\item[\code{phi}] A vector for the autoregressive (AR) parameters for AR part.

\item[\code{theta}] A vector for the moving average (MA) parameters for MA part.

\item[\code{Delta}] The parameter vector contains the changepoint parameters for time series mean function.

\item[\code{CpLoc}] A vector contains the changepoint locations range from \eqn{1\leq\tau\leq T_{s}}{}.

\item[\code{seed}] The random seed for simulation reproducibility.
\end{ldescription}
\end{Arguments}
%
\begin{Details}
The simulated time series \eqn{Z_{t}, t=1,\ldots,T_{s}}{} is from a class of models,
\deqn{Z_{t}=\mu_{t}+e_{t}.}{}
\begin{itemize}

\item{} Time series observations are IID \\{} \eqn{\mu_{t}}{} is a constant
and \eqn{e_{t}}{}'s are independent and identically distributed as \eqn{N(0,\sigma)}{}.
\item{} ARMA(p,q) model with constant mean \\{} \eqn{\mu_{t}}{} is a constant,
and \eqn{e_{t}}{} follows an ARMA(p,q) process.
\item{} ARMA(p,q) model with seasonality \\{} \eqn{\mu_{t}}{} is not a constant,
\deqn{\mu_{t} = ASin(\frac{2\pi t}{S})+BSin(\frac{2\pi t}{S}),}{}
and \eqn{e_{t}}{} follows an ARMA(p,q) process.
\item{} ARMA(p,q) model with seasonality and trend \\{} \eqn{\mu_{t}}{} is not a constant,
\deqn{\mu_{t} = ASin(\frac{2\pi t}{S})+BSin(\frac{2\pi t}{S}) + \alpha t,}{}
and \eqn{e_{t}}{} follows an ARMA(p,q) process.
The changepoint effects could be introduced through the \eqn{\mu_{t}}{} as
\deqn{\mu_{t} = \Delta_{1}I_{t>\tau_{1}} + \ldots + \Delta_{m}I_{t>\tau_{m}},}{}
where \eqn{1\leq\tau_{1}<\ldots<\tau_{m}\leq T_{s}}{} are the changepoint locations and
\eqn{\Delta_{1},\ldots,\Delta_{m}}{} are the changepoint parameter that need to be estimated.

\end{itemize}

\end{Details}
%
\begin{Value}
The simulated time series with attributes:
\begin{ldescription}
\item[\code{\code{Z}}] The simulated time series.
\item[\code{Attributes}] 
\begin{itemize}

\item{} \code{DesignX} The covariates include all changepoint indicators.
\item{} \code{mu} A vector includes the mean values for simulated time series sequences.
\item{} \code{theta} The true parameter vector (including changepoint effects).
\item{} \code{CpEff} The true changepoint magnitudes.
\item{} \code{CpLoc} The true changepoint locations where we introduce the mean changing effects.
\item{} \code{arEff} A vector givies the AR coefficients.
\item{} \code{maEff} A vector givies the MA coefficients.
\item{} \code{seed} The random seed used for this simulation.

\end{itemize}


\end{ldescription}
\end{Value}
%
\begin{Examples}
\begin{ExampleCode}
##### M1: Time series observations are IID
Ts = 1000
betaT = c(0.5) # intercept
XMatT = matrix(1, nrow=Ts, ncol=1)
colnames(XMatT) = "intercept"
sigmaT = 1
DeltaT = c(2, -2)
Cp.prop = c(1/4, 3/4)
CpLocT = floor(Ts*Cp.prop)

myts = ts.sim(beta=betaT, XMat=XMatT, sigma=sigmaT,
              Delta=DeltaT, CpLoc=CpLocT, seed=1234)
TsPlotCheck(Y=myts, tau=CpLocT)

##### M2: ARMA(2,1) model with constant mean
Ts = 1000
betaT = c(0.5) # intercept
XMatT = matrix(1, nrow=Ts, ncol=1)
colnames(XMatT) = "intercept"
sigmaT = 1
phiT = c(0.5, -0.5)
thetaT = c(0.8)
DeltaT = c(2, -2)
Cp.prop = c(1/4, 3/4)
CpLocT = floor(Ts*Cp.prop)

myts = ts.sim(beta=betaT, XMat=XMatT, sigma=sigmaT,
              phi=phiT, theta=thetaT, Delta=DeltaT, CpLoc=CpLocT, seed=1234)
TsPlotCheck(Y=myts, tau=CpLocT)

##### M3: ARMA(2,1) model with seasonality
Ts = 1000
betaT = c(0.5, -0.5, 0.3) # intercept, B, D
period = 30
XMatT = cbind(rep(1, Ts), cos(2*pi*(1:Ts)/period), sin(2*pi*(1:Ts)/period))
colnames(XMatT) = c("intercept", "Bvalue", "DValue")
sigmaT = 1
phiT = c(0.5, -0.5)
thetaT = c(0.8)
DeltaT = c(2, -2)
Cp.prop = c(1/4, 3/4)
CpLocT = floor(Ts*Cp.prop)

myts = ts.sim(beta=betaT, XMat=XMatT, sigma=sigmaT,
              phi=phiT, theta=thetaT, Delta=DeltaT, CpLoc=CpLocT, seed=1234)
TsPlotCheck(Y=myts, tau=CpLocT)


##### M4: ARMA(2,1) model with seasonality and trend
# scaled trend if large number of sample size
Ts = 1000
betaT = c(0.5, -0.5, 0.3, 0.01) # intercept, B, D, alpha
period = 30
XMatT = cbind(rep(1, Ts), cos(2*pi*(1:Ts)/period), sin(2*pi*(1:Ts)/period), 1:Ts)
colnames(XMatT) = c("intercept", "Bvalue", "DValue", "trend")
sigmaT = 1
phiT = c(0.5, -0.5)
thetaT = c(0.8)
DeltaT = c(2, -2)
Cp.prop = c(1/4, 3/4)
CpLocT = floor(Ts*Cp.prop)

myts = ts.sim(beta=betaT, XMat=XMatT, sigma=sigmaT,
              phi=phiT, theta=thetaT, Delta=DeltaT, CpLoc=CpLocT, seed=1234)
TsPlotCheck(Y=myts, tau=CpLocT)
\end{ExampleCode}
\end{Examples}
\inputencoding{utf8}
\HeaderA{TsPlotCheck}{Plot the simulated time series}{TsPlotCheck}
%
\begin{Description}
This is a function to plot the simulated time series with segmentation visualization
by provided changepoint locations.
\end{Description}
%
\begin{Usage}
\begin{verbatim}
TsPlotCheck(
  X = NULL,
  Xat = NULL,
  Y,
  tau = NULL,
  mu = NULL,
  XLAB = NULL,
  YLAB = NULL
)
\end{verbatim}
\end{Usage}
%
\begin{Arguments}
\begin{ldescription}
\item[\code{X}] The time series time index, which could be specified as years, months, days, or others.
The default value is NULL and the vector from 1 to the time series length will be applied.

\item[\code{Xat}] The values from \code{X} that will be used as the X axis tick marks.

\item[\code{Y}] The time series data.

\item[\code{tau}] The provided changepoint locations.

\item[\code{mu}] The provided meam values for each time \eqn{t}{}.

\item[\code{XLAB}] A descriptive label for X axis.

\item[\code{YLAB}] A descriptive label for Y axis.
\end{ldescription}
\end{Arguments}
%
\begin{Examples}
\begin{ExampleCode}
Ts = 1000
betaT = c(0.5) # intercept
XMatT = matrix(1, nrow=Ts, ncol=1)
colnames(XMatT) = "intercept"
sigmaT = 1
DeltaT = c(2, -2)
Cp.prop = c(1/4, 3/4)
CpLocT = floor(Ts*Cp.prop)

myts = ts.sim(beta=betaT, XMat=XMatT, sigma=sigmaT, Delta=DeltaT, CpLoc=CpLocT, seed=1234)
TsPlotCheck(X=1:Ts, Xat=seq(from=1, to=Ts, length=10), Y=myts, tau=CpLocT)
\end{ExampleCode}
\end{Examples}
\inputencoding{utf8}
\HeaderA{uniformcrossover}{Uniform crossover to produce offsprings}{uniformcrossover}
%
\begin{Description}
In uniform crossover, typically, each bit is chosen from either parent with
equal probability. Other mixing ratios are sometimes used, resulting in
offspring which inherit more genetic information from one parent than the
other. In a uniform crossover, we don’t divide the chromosome into segments,
rather we treat each gene separately. In this, we essentially flip a coin
for each chromosome to decide whether or not it will be included in the
off-spring. If model order selection is requested, each child's model order
has the equal probability (0.5) from \code{dad} and \code{mom}.
\end{Description}
%
\begin{Usage}
\begin{verbatim}
uniformcrossover(mom, dad, prange, minDist, lmax, N)
\end{verbatim}
\end{Usage}
%
\begin{Arguments}
\begin{ldescription}
\item[\code{mom}] Among two selected individuals, \code{mom} represents the selected
chromosome representation with lower fitness function value.

\item[\code{dad}] Among two selected individuals, \code{dad} represents the selected
chromosome representation with larger fitness function value.

\item[\code{prange}] Default value is \code{NULL} for only changepoint detection. If
\code{prange} is specified as a list object, which contains the range of
each model order parameters for order selection (integers). The number of
order parameters must be equal to the length of \code{prange}.

\item[\code{minDist}] The required minimum distance between two adjacent changepoints.

\item[\code{lmax}] The user specified maximum number of changepoints, by default,
as \code{N/2 - 1}.

\item[\code{N}] The length of time series.
\end{ldescription}
\end{Arguments}
%
\begin{Value}
The child chromosome that produced from \code{mom} and \code{dad} for
next generation.
\end{Value}
\printindex{}
\end{document}
